% !TEX root = ../aa_SoM_Chapters.tex

% ö = \%c3\%b6
% ä = \%C3\%A4

\xdef\sectiontitle{\Isectiontitle} %aiheuttama
\TOCrefs

%%%%%%%%%%%%%%%
\begin{frame}[label=1_11_a]%[label=1_2_TOC1]
\frametitle{\texorpdfstring{\culine[\sectiontitleulinecolor]{\sectiontitle}}{\sectiontitle} }
\label{\TOCname}

\vspace{0.5cm}
% can we simply override the appearing of items when needed with <1-> etc.
\begin{itemize}[leftmargin=0.5cm,itemsep=2mm]
    \item[1.1]<1-> \hyperlink{1_1_a}{\IaSubsectiontitle}
    \item[1.2]<1-> \hyperlink{1_2_a}{\IbSubsectiontitle}	
    \item[1.3]<1-> \hyperlink{1_3_a}{\IcSubsectiontitle}	
    \item[1.4]<1-> \hyperlink{1_4_a}{\IdSubsectiontitle}
    \item[1.5]<1-> \hyperlink{1_5_a}{\IeSubsectiontitle}
    \item[1.6]<1-> \hyperlink{1_6_a}{\IfSubsectiontitle}
    \item[1.7]<1-> \hyperlink{1_7_a}{\IgSubsectiontitle}
\end{itemize}

\vspace{-1cm}
\begin{itemize}[leftmargin=8.5cm,itemsep=2mm]
    \item[1.8]<1-> \hyperlink{1_8_a}{\IhSubsectiontitle}
	\item[1.9]<1-> \hyperlink{1_9_a}{\IiSubsectiontitle}
	\item[1.10]<1-> \hyperlink{1_10_a}{\IjSubsectiontitle}
	\item[1.11]<1-> \hyperlink{1_11_a}{\CDAlert[blue]{\IkSubsectiontitle}}
\end{itemize}

%\endofslide 
\end{frame}
%%%%%%%%%%%%%%%

\xdef\sectiontitle{\IkSubsectiontitle} 

%%%%%%%%%%%%%%%
\begin{frame}%[label=1_11_a]
\frametitle{\texorpdfstring{\culine[\sectiontitleulinecolor]{\sectiontitle}}{\sectiontitle} }
%\framesubtitle{Phonons}

\begin{itemize}[itemsep=1.5mm]
	\item There is a lower limit to phonon wavelengths 
	\vspace{-1mm}\implies{In practice: \qquad $\frac{1}{2} \appear{\lambda_{\min}}  \equal \appear{a=$} \appear{the lattice constant} } 
	
	\item A minimum wavelength indicates maximum frequency \appear{and maximum energy $E_{\max }$ for the phonons}
\end{itemize}

% 6.5 cm = midway, 10 cm = 3/4 
%\vspace{2.5mm}
\begin{minipage}{8.5cm}
\begin{itemize}[itemsep=1.5mm]
	\item The \CDAlert[fuchsia]{speed of phonons} is not the speed of light\appear{, but rather the \CDAlert[fuchsia]{speed of sound $v_{s}$} in the material}

\item Phonons can oscillate along two \CDAlert[red]{transverse} directions\appear{, but also \CDAlert[blue]{longitudinal} vibrations are allowed}

\item With equal speeds for both transverse and longitudinal oscillations\appear{, the phononic density of states must be multiplied by $3 / 2$}
\end{itemize}
\end{minipage}


\xdef\loc{south east}
\begin{tikzpicture}[remember picture,overlay] % anchor=south
    \node (A) [yshift=-0*\figYshift,xshift=\figXshift, anchor=\loc] at (current page.\loc) {\includegraphics[page=1,height=6.7cm]{Figures_booklet/SOM_9_Statistical_mechanics_figures/SOM_Figure_30_cropped.png}}; %scale=\figscale. % 8cm max hei
\end{tikzpicture}

\endofslide 
%\endofsection
\end{frame}
%%%%%%%%%%%%%%%

%%%%%%%%%%%%%%%
\begin{frame}
\frametitle{\texorpdfstring{\culine[\sectiontitleulinecolor]{\sectiontitle}}{\sectiontitle} }

\begin{itemize}[itemsep=1.55mm]
	\item Now we can write the formula for the energy of a phonon system:
\[
U_{\text {phonon }}  \equal \appear{\nint_{0}^{E_{\max}}} \appear{E} \appear{\mathbb{N}(E)} \appear{D(E)} \appear{d E}
 \equal 
 \appear{\nint_{0}^{E_{\max }} E} \frac{1}{e^{E / k_{B} T}-1} \frac{12 \pi V}{h^{3} v_{S}{ }^{3}} \appear{E^{2}} \appear{d E}
\]

\item To perform the integral\appear{, we need an \Alert[red]{estimate for the maximum}}
% 6.5 cm = midway, 10 cm = 3/4 
\vspace{2.5mm}
\end{itemize}

\begin{minipage}{7.3cm}
\begin{itemize}[itemsep=1.55mm]

	 \item[] \appear{\Alert[red]{energy} $E_{\max }$}\appear{, and link it to the smallest wavelength} \appear{$2 a$}

\item Consider the atoms in a \Alert[blue]{cubic volume of a crystalline solid} 

\item Number of atoms on one edge is $N^{1 / 3}$\appear{, and spacing between atoms is $L / N^{1 / 3}$} 
\item[] \vspace{-2.2mm}\implies{ Solid supports phonons, with the smallest phonon wavelength being $2 L / N^{1 / 3}$}

\end{itemize}
\end{minipage}


\xdef\loc{south east}
\begin{tikzpicture}[remember picture,overlay] % anchor=south
    \node (A) [yshift=-0.25*\figYshift,xshift=\figXshift, anchor=\loc] at (current page.\loc) {\includegraphics[page=1,height=5.35cm]{Figures_booklet/SOM_9_Statistical_mechanics_figures/SOM_Figure_31.png}}; %scale=\figscale. % 8cm max hei
\end{tikzpicture}

\endofslide 
%\endofsection
\end{frame}
%%%%%%%%%%%%%%%

%%%%%%%%%%%%%%%%
%\begin{frame}
%\frametitle{\texorpdfstring{\culine[\sectiontitleulinecolor]{\sectiontitle}}{\sectiontitle} }
%
%\begin{itemize}
%	\item Minimum wavelength leads to a maximum wave number $k=\frac{2 \pi}{\lambda}  \equal  \frac{\pi N^{1 / 3}}{L }$ 
%	
%	\item Note, $\mathbf{k}$ is basically a vector (k=$|\mathbf{k}|$)\appear{, it may point at any direction}\appear{, allowed phononic states form a sphere!}
%% 6.5 cm = midway, 10 cm = 3/4 
%\item[] \begin{minipage}{8.5cm}
%\item We consider a cube in $\mathbf{k}$-space, with edge of $\pi N^{1 / 3} / L$, and we set the volume of the cube to be equal to $1 / 8$ of a sphere with radius $k_{\max }$
%\[
%\begin{aligned}
% \textrm{((} \frac{\pi N^{1 / 3}}{L} \textrm{))} ^{3} &= \frac{1}{8} \textrm{((} \frac{4 \pi}{3} k_{\max }^{3} \textrm{))}  \\
%\Rightarrow \qquad k_{\max } &= \textrm{((} \frac{3}{4 \pi} 8 \textrm{((} \frac{\pi N^{1 / 3}}{L} \textrm{))} ^{3} \textrm{))} ^{1 / 3} \\
%&= \textrm{((} \frac{6 \pi^{2} N}{L^{3}} \textrm{))} ^{1 / 3}= \textrm{((} \frac{6 \pi^{2} N}{V} \textrm{))} ^{1 / 3}
%\end{aligned}
%\]
%\end{minipage}
%\vspace{3mm}
%\end{itemize}
%
%
%\xdef\loc{south east}
%\begin{tikzpicture}[remember picture,overlay] % anchor=south
%    \node (A) [yshift=-0.25*\figXshift,xshift=\figXshift, anchor=\loc] at (current page.\loc) {\includegraphics[page=1,height=7cm]{Figures_booklet/SOM_9_Statistical_mechanics_figures/SOM_Figure_32.png}}; %scale=\figscale. % 8cm max hei
%\end{tikzpicture}
%
%
%\endofslide 
%%\endofsection
%\end{frame}
%%%%%%%%%%%%%%%%
%

%%%%%%%%%%%%%%%
\begin{frame}
\frametitle{\texorpdfstring{\culine[\sectiontitleulinecolor]{\sectiontitle}}{\sectiontitle} }

\begin{itemize}[itemsep=1.7mm] \fontsize{11.5}{13}\selectfont %\small
%	\item Minimum wavelength leads to a maximum wave number $k=\frac{2 \pi}{\lambda}  \equal  \frac{\pi N^{1 / 3}}{L }$ 
	\item \CDAlert[blue]{Due to the periodicity} of the solid\appear{, the allowed \Alert[blue]{wave vector $\mathbf{k}$ values are discretized}:} \appear{a single phononic state occupies a volume of  $( 2 \pi /L )^{3}$} \appear{(in the so-called $\mathbf{k}$-space}\appear{/state space)}
\end{itemize}
\vspace{1.5mm}
	 
% 6.5 cm = midway, 10 cm = 3/4 
\begin{minipage}{8cm} %\small
\begin{itemize}[itemsep=1.7mm] \fontsize{11.5}{13}\selectfont %\small

\item Note, $\mathbf{k}$ is basically a vector \appear{(k=$|\mathbf{k}|$)}\appear{, it may point at any direction}\appear{, allowed phononic states form a sphere!}

\item \CDAlert[red]{The system has in total $N$ states}\appear{, which must be equal to the product of the density} \appear{and the volume of states} \appear{(in $\mathbf{k}$-space):}
\[
\begin{aligned}
N & \equal  \appear{\bigg( \Frac{L}{2\pi} \bigg)^{3}} \appear{\bigg( \Frac{4 \pi}{3} k_{\max }^{3} \bigg)}  \\
\appear{\Rightarrow} \qquad \appear{k_{\max }} & \equal  \appear{\bigg( \frac{8\pi^3}{L^3}} \frac{3N}{4\pi} \appear{\bigg) ^{1 / 3}} \\
& \equal  \appear{\bigg(} \frac{6 \pi^{2} N}{L^{3}} \appear{\bigg) ^{1 / 3}}  \equal  \appear{\bigg(} \frac{6 \pi^{2} N}{V} \appear{\bigg) ^{1 / 3}}
\end{aligned}
\]
\end{itemize}
\end{minipage}


\xdef\loc{south east}
\begin{tikzpicture}[remember picture,overlay] % anchor=south
    \node (A) [yshift=-0.25*\figXshift,xshift=\figXshift, anchor=\loc] at (current page.\loc) {\includegraphics[page=1,height=7cm]{Figures_booklet/SOM_9_Statistical_mechanics_figures/SOM_Figure_32.png}}; %scale=\figscale. % 8cm max hei
\end{tikzpicture}


\endofslide 
%\endofsection
\end{frame}
%%%%%%%%%%%%%%%

%%%%%%%%%%%%%%%
\begin{frame}
\frametitle{\texorpdfstring{\culine[\sectiontitleulinecolor]{\sectiontitle}}{\sectiontitle} }

\begin{itemize}
	\item Having derived the maximum wave number\appear{, we can now estimate maximum energy of the phonons:}
\[
E_{\max }  \equal  \frac{h v_{s} k_{\max }}{2 \pi}  \equal  \frac{h v_{s}}{2 \pi} \appear{\textrm{((}} \frac{6 \pi^{2} N}{V} \appear{\textrm{))} ^{1 / 3}}  \equal  \appear{h v_{s}} \appear{\textrm{((}} \frac{3 N}{4 \pi V} \appear{\textrm{))} ^{1 / 3}} \appear{\,, \qquad E=\hbar \omega ~\textrm{and}~\omega = v_s k}
\]

\item Here, we have \CDAlert[red]{assumed} \appear{that the speed of the propagating waves in the solid is the same as the speed of sound}\appear{, verified experimentally}\appear{, and that \CDAlert[red]{all phonons move with same speed}} 

\item Setting $E_{\max }=k_{B} T_{D}$ \appear{we get the maximum energy expressed in units of temperature}\appear{, with so-called \CDAlert[blue]{Debye temperature} }\appear{, which obviously depends on properties of the given material}\appear{, such as $N / V$ and $v_{s}$:}
\[
T_{D}  \equal \frac{E_{\max }}{k_{B}}  \equal \frac{h v_{S}}{k_{B}} \appear{\bigg(} \frac{3 N}{4 \pi V} \appear{\bigg) ^{1 / 3}}
\]
\end{itemize}

\endofslide 
%\endofsection
\end{frame}
%%%%%%%%%%%%%%%

%%%%%%%%%%%%%%%
\begin{frame}
\frametitle{\texorpdfstring{\culine[\sectiontitleulinecolor]{\sectiontitle}}{\sectiontitle} }

\begin{itemize}
	\item With the maximum energy derived\appear{, we can return to the original integral for the total energy of the system:}
	\appear{$$
	U_{\text {phonon}}  \equal   \appear{\nint_{0}^{k_{B} T_{D}}} \appear{E} \fracc{1}{e^{E / k_{B} T}-1} \fracc{9 N}{ \textrm{((} k_{B} T_{D} \textrm{))} ^{3}} \appear{E^{2}} \appear{d E}  \equal \appear{9 k_{B} N} \frac{T^{4}}{T_{D}^{3}} \appear{\nint_{0}^{T_{D} / T}} \frac{x^{3}}{e^{x}-1} \appear{d x}
	$$
	\appear{where we have changed the variable $E$} \appear{into $x$} $\appear{(x} \equiv \appear{E/k_{B} T)}$}
\item For a mole of material $N=N_{A}$\appear{, and since $k_{B} N_{A}=R$}\appear{, we obtain finally:}
\[
U_{\text{phonon,\,per mole}}  \equal \appear{9 R} \frac{T^{4}}{T_{D}^{3}} \appear{\nint_{0}^{T_{D} / T}} \frac{x^{3}}{e^{x}-1} \appear{d x} \qquad
\]

\item Unfortunately, the \CDAlert[red]{integral cannot be calculated analytically}\appear{, we need to rely on numerical integration}\appear{, and compare with experimental data}
\end{itemize}


\endofslide 
%\endofsection
\end{frame}
%%%%%%%%%%%%%%%

%%%%%%%%%%%%%%%
\begin{frame}
\frametitle{\texorpdfstring{\culine[\sectiontitleulinecolor]{\sectiontitle}}{\sectiontitle} }

\begin{itemize}
	\item The \CDAlert[fuchsia]{Debye model} described above\appear{, \Alert[fuchsia]{fits well with experimental data on heat capacities}}\appear{, as is seen in the below Figure}
	
	\item \appear{The \CDAlert[red]{derivation was quick\;\&\;dirty}}\appear{, rigorous one will be given in our introductory course on solid-state physics!}
\end{itemize}

\xdef\loc{south}
\begin{tikzpicture}[remember picture,overlay] % anchor=south
    \node (A) [yshift=-0.25*\figYshift,xshift=\figXshift, anchor=\loc] at (current page.\loc) {\includegraphics[page=1,width=14cm]{Figures_booklet/SOM_9_Statistical_mechanics_figures/SOM_Figure_33.png}}; %scale=\figscale. % 8cm max hei
\end{tikzpicture}


%\endofslide 
\endofsection
\end{frame}
%%%%%%%%%%%%%%%
